\documentclass[9pt]{beamer}
\usetheme{metropolis}
\usepackage{amsmath}
\usepackage{amssymb}
\usepackage{graphicx}
\usepackage{booktabs}
\usepackage{xcolor}
\usepackage{listings}

% Define color palette matching our project design system
\definecolor{primary}{HTML}{1abc9c}    % Deep Teal
\definecolor{secondary}{HTML}{3498db}  % Astro Blue
\definecolor{darkbg}{HTML}{2c3e50}     % Slate Blue
\definecolor{accent}{HTML}{e67e22}     % Orange

\setbeamercolor{palette primary}{bg=darkbg, fg=white}
\setbeamercolor{palette secondary}{bg=primary, fg=white}
\setbeamercolor{palette tertiary}{bg=secondary, fg=white}
\setbeamercolor{structure}{fg=primary}
\setbeamercolor{alerted text}{fg=accent}

\title{Stellar Trace: Host Profiling \& Machine Learning}
\subtitle{Meet 1: Observational Selection \& Non-Parametric Statistics (Tasks 7--9)}
\author{Astrophysics Mentorship Program}
\date{\today}

\begin{document}

\maketitle

% Slide 2: Meet Series Roadmap
\begin{frame}{Meet Series Roadmap}
  \begin{block}{Phase 2 Curriculum Structure}
    Stellar Trace Phase 2 transitions from the volume-complete universe model of Milestone 1 to real observational data ingestion, classifier training, and physical model fitting:
  \end{block}
  \begin{description}
    \item[Meet 1 (Today):] \textbf{Telescopes \& Statistics (Tasks 7--9)}
      \begin{itemize}
        \item Task 7: Observational Selection Effects \& Malmquist Bias.
        \item Task 8: Ingestion, Coordinate Cross-Matching, and SQL WAF Bypass.
        \item Task 9: 1D vs. 2D Kolmogorov-Smirnov statistical testing.
      \end{itemize}
    \item[Meet 2:] \textbf{Machine Learning \& Data Leakage (Tasks 10--12)}
      \begin{itemize}
        \item Oversampling under strict train-test separation; MLP Classifiers.
      \end{itemize}
    \item[Meet 3:] \textbf{Rate Fitting \& Physics-Informed nets (Tasks 13--15)}
      \begin{itemize}
        \item Bayesian MCMC (Metropolis-Hastings), PINNs, and DTD Convolution.
      \end{itemize}
  \end{description}
\end{frame}

% Slide 3: Meet 1 Objectives
\begin{frame}{Meet 1 Objectives}
  \begin{itemize}
    \item In \textbf{Milestone 1}, you generated a volume-complete synthetic universe based on galaxy mass functions and star-forming correlations.
    \item \alert{The Problem:} Real telescope catalogs are not volume-complete. They are limited by the telescope's mirror size, detector sensitivity, and atmospheric dust.
    \item \textbf{Today's Goals:}
      \begin{enumerate}
        \item Mathematically apply telescope sensitivity thresholds (\textbf{Malmquist Bias}) to our mock universe.
        \item Ingest real Core-Collapse and Type Ia supernova host catalog data.
        \item Write coordinate cross-matching algorithms to link supernovae to their hosts.
        \item Query the SDSS spectroscopic database server and bypass firewall restrictions.
        \item Run 1D and 2D non-parametric statistics to prove that supernovae live in distinct environments.
      \end{enumerate}
  \end{itemize}
\end{frame}

% Slide 4: Section 1 Divider
\begin{frame}[plain]
  \vfill
  \centering
  \begin{beamercolorbox}[sep=8pt,center,shadow=true,rounded=true]{palette primary}
    \usebeamerfont{title}\textbf{Section 1: Observational Selection \& Malmquist Bias (Task 7)}
  \end{beamercolorbox}
  \vfill
\end{frame}

% Slide 5: Volume-Complete vs. Flux-Limited Surveys
\begin{frame}{Volume-Complete vs. Flux-Limited Surveys}
  To understand real telescope data, we must differentiate between two types of astronomical surveys:
  \begin{itemize}
    \item \textbf{Volume-Complete Survey:}
      Contains every single galaxy within a specified physical volume $V$ (comoving sphere), regardless of how bright they appear. Extremely difficult to achieve because small, faint dwarf galaxies are missed.
    \item \textbf{Flux-Limited Survey:}
      Only includes objects whose flux on the sky is higher than the telescope's detector sensitivity threshold:
      $$f > f_{\text{limit}}$$
    \item \textbf{SDSS Spectroscopic Limit:}
      The Sloan Digital Sky Survey spectroscopic catalog is complete down to a limiting apparent magnitude of $m_r < 23.5$ in the $r$-band. Any galaxy fainter than this is ignored.
  \end{itemize}
\end{frame}

% Slide 6: The Physics of Malmquist Bias
\begin{frame}{The Physics of Malmquist Bias}
  \begin{block}{Light Propagation}
    The flux of light received on Earth scales inversely with the square of the luminosity distance $D_L$:
    $$f = \frac{L}{4\pi D_L^2}$$
  \end{block}
  \begin{itemize}
    \item For a telescope with a fixed flux limit $f_{\text{limit}}$, the minimum detectable luminosity increases quadratically with distance:
      $$L_{\text{min}}(z) = 4\pi D_L(z)^2 f_{\text{limit}}$$
    \item At low redshifts (close to us), we can detect both massive, bright galaxies and tiny, faint dwarf galaxies.
    \item At high redshifts (far away), dwarf galaxies fall below the detection limit and disappear from our survey.
    \item \alert{Malmquist Bias:} The systematic selection of high-luminosity objects at large distances. If uncorrected, we would falsely conclude that galaxies in the distant universe are systematically more massive and active than today.
  \end{itemize}
\end{frame}

% Slide 7: Mathematical Model Derivation: Step 1 & 2
\begin{frame}{Mathematical Model Derivation: Step 1 \& 2}
  To apply the survey limit ($m_r < 23.5$) to our mock galaxies, we must convert their simulated properties ($M_\star, \text{sSFR}, z$) into $r$-band apparent magnitude ($m_r$):
  \begin{enumerate}
    \item \textbf{Step 1: Specific Star Formation Rate (sSFR) to rest-frame $(g-r)$ color:}
      Young, massive stars are hot and blue; old, low-mass stars are cool and red. Active star-forming galaxies (high sSFR) are blue, and quenched galaxies are red:
      $$(g-r) = \text{clip}(0.5 - 0.25 \cdot (\text{sSFR} + 10), -0.1, 1.0) + \eta$$
      where $\eta \sim \mathcal{N}(0, 0.08)$ represents intrinsic astrophysical scatter.
    \item \textbf{Step 2: Bell \& de Jong (2001) Mass-to-Light ($M/L$) ratio in $g$-band:}
      Stellar population synthesis models relate color directly to $M/L$ in the $g$-band:
      $$\log_{10}(M/L)_g = 1.5 \cdot (g-r) - 0.7$$
  \end{enumerate}
\end{frame}

% Slide 7b: Detailed Derivation: Translation from g-band to r-band
\begin{frame}{Detailed Derivation: Translating $M/L$ from $g$-band to $r$-band}
  How do we find the Mass-to-Light ratio in the $r$-band, $(M/L)_r$?
  \begin{itemize}
    \item The rest-frame color $(g-r)$ is defined as the difference in absolute magnitudes:
      $$(g-r) = M_g - M_r$$
    \item Substituting absolute magnitudes in Solar units, this relates to the luminosity ratio:
      $$(g-r) = -2.5 \log_{10} \left( \frac{L_g / L_{\odot, g}}{L_r / L_{\odot, r}} \right) \implies \log_{10} \left( \frac{\mathcal{L}_g}{\mathcal{L}_r} \right) = -\frac{g-r}{2.5}$$
    \item Since the physical stellar mass $M_\star$ of a galaxy is identical in both bands:
      $$\frac{(M/L)_r}{(M/L)_g} = \frac{M_\star / \mathcal{L}_r}{M_\star / \mathcal{L}_g} = \frac{\mathcal{L}_g}{\mathcal{L}_r}$$
    \item Taking the logarithm of both sides gives the translation formula:
      $$\log_{10}(M/L)_r = \log_{10}(M/L)_g - \frac{g-r}{2.5}$$
  \end{itemize}
\end{frame}

% Slide 8: Mathematical Model Derivation: Step 3
\begin{frame}{Mathematical Model Derivation: Step 3}
  \begin{enumerate}
    \setcounter{enumi}{2}
    \item \textbf{Step 3: Absolute Magnitude ($M_r$):}
      Absolute magnitude represents the galaxy's intrinsic luminosity. We calculate $M_r$ relative to the Solar absolute magnitude in the $r$-band ($M_{\odot, r} = 4.65$):
      $$M_r = M_{\odot, r} - 2.5 \log_{10} \left( \frac{L_r}{L_{\odot, r}} \right)$$
    \item By definition of the mass-to-light ratio, the luminosity in Solar units is:
      $$\frac{L_r}{L_{\odot, r}} = \frac{M_\star}{(M/L)_r}$$
    \item Substituting this yields our absolute magnitude formula:
      $$M_r = 4.65 - 2.5 \cdot \left( \log_{10} M_\star - \log_{10}(M/L)_r \right)$$
      where $M_\star$ is the stellar mass in Solar units.
  \end{enumerate}
\end{frame}

% Slide 8b: Mathematical Model Derivation: Step 4
\begin{frame}{Mathematical Model Derivation: Step 4}
  \begin{enumerate}
    \setcounter{enumi}{3}
    \item \textbf{Step 4: Cosmology and Luminosity Distance ($D_L(z)$):}
      In a flat $\Lambda$CDM cosmology ($H_0 = 70.0 \, \text{km/s/Mpc}, \Omega_{m,0} = 0.3$), the comoving distance is integrated:
      $$D_C(z) = \frac{c}{H_0} \int_0^z \frac{dz'}{\sqrt{\Omega_{m,0}(1+z')^3 + \Omega_{\Lambda,0}}}$$
    \item The luminosity distance is:
      $$D_L(z) = (1+z) \cdot D_C(z)$$
      To speed up computations, we pre-compute a grid of $(z, D_L)$ and interpolate using \texttt{scipy.interpolate.interp1d}.
  \end{enumerate}
\end{frame}

% Slide 9: Detailed Derivation: Distance Modulus & Apparent Magnitude
\begin{frame}{Detailed Derivation: Distance Modulus \& Apparent Magnitude}
  \begin{itemize}
    \item Absolute magnitude $M_r$ is defined at a standard distance of exactly $10\text{ pc}$. Since flux scales as $f \propto 1/d^2$:
      $$\frac{f(D_L)}{f(10\text{ pc})} = \left( \frac{10\text{ pc}}{D_L} \right)^2 \implies m_r - M_r = 5 \log_{10} \left( \frac{D_L}{10\text{ pc}} \right)$$
    \item To use $D_L(z)$ in Megaparsecs ($\text{Mpc}$), where $1\text{ Mpc} = 10^6\text{ pc}$, we convert:
      $$\frac{D_L(\text{in pc})}{10\text{ pc}} = \frac{D_L(\text{in Mpc}) \times 10^6}{10} = D_L(\text{in Mpc}) \times 10^5$$
    \item Substitute into the distance modulus equation:
      $$m_r - M_r = 5 \log_{10} \left( D_L(\text{Mpc}) \times 10^5 \right) = 5 \log_{10} D_L(\text{Mpc}) + 25$$
    \item This gives the final apparent magnitude formula:
      $$m_r = M_r + 5 \log_{10} D_L(\text{Mpc}) + 25$$
  \end{itemize}
\end{frame}

% Slide 9b: Worked Example: Malmquist Bias in Action
\begin{frame}{Worked Example: Malmquist Bias in Action}
  Let's compute apparent magnitudes for a galaxy with $M_\star = 10^{10} M_\odot$ ($\log_{10} M_\star = 10.0$) and color $(g-r) = 0.6$:
  \begin{itemize}
    \item \textbf{1. Mass-to-light ratio:} $\log_{10}(M/L)_g = 1.5(0.6) - 0.7 = 0.2$
    \item \textbf{2. translation to r-band:} $\log_{10}(M/L)_r = 0.2 - 0.6/2.5 = -0.04$
    \item \textbf{3. Absolute magnitude:} $M_r = 4.65 - 2.5(10.0 - (-0.04)) = -20.45$
  \end{itemize}
  Now let's compare two distances:
  \begin{enumerate}
    \item \textbf{Close proximity ($z \approx 0.08, D_L = 360\text{ Mpc}$):}
      $$m_r = -20.45 + 5\log_{10}(360) + 25 = 17.33 \quad \implies \quad \textbf{Detected } (17.33 < 23.5)$$
    \item \textbf{Distant universe ($z \approx 0.95, D_L = 6200\text{ Mpc}$):}
      $$m_r = -20.45 + 5\log_{10}(6200) + 25 = 23.51 \quad \implies \quad \textbf{Masked Out } (23.51 \ge 23.5)$$
  \end{enumerate}
  This proves why distant dwarf galaxies vanish from our survey!
\end{frame}

% Slide 9c: Mathematical Model Derivation: Step 6
\begin{frame}{Mathematical Model Derivation: Step 6}
  \begin{enumerate}
    \setcounter{enumi}{5}
    \item \textbf{Step 6: Selection Mask:}
      We apply the spectroscopic detection filter:
      $$\text{detected\_mask} = m_r < 23.5$$
      Galaxies with $m_r \ge 23.5$ are masked out, transforming our volume-complete catalog into a flux-limited observational catalog.
  \end{enumerate}
  \begin{alertblock}{Computational Consequence}
    Applying this mask reduces our mock sample size from $\sim 100,000$ to $\sim 45,000$ galaxies, accurately mimicking real spectroscopic selection effects.
  \end{alertblock}
\end{frame}

% Slide 10: Section 2 Divider
\begin{frame}[plain]
  \vfill
  \centering
  \begin{beamercolorbox}[sep=8pt,center,shadow=true,rounded=true]{palette primary}
    \usebeamerfont{title}\textbf{Section 2: Supernova Progenitors \& Data Ingestion (Task 8)}
  \end{beamercolorbox}
  \vfill
\end{frame}

% Slide 11: Supernova Progenitors: Core-Collapse SNe
\begin{frame}{Supernova Progenitors: Core-Collapse SNe}
  \begin{itemize}
    \item \textbf{Progenitors:} Massive stars ($M \ge 8 M_\odot$, spectral classes O and B).
    \item \textbf{Lifespans:} Extremely short ($\tau \le 50\text{ Myr}$), meaning they explode almost immediately after their parent molecular cloud collapses.
    \item \textbf{Host Environment:}
      Because massive stars die before they can migrate out of their birth environments, Core-Collapse SNe directly trace instantaneous star formation. They reside exclusively in active star-forming galaxies (spirals and starbursts).
    \item \textbf{Data source:} Confirmed observational catalog table from \textit{Sharma et al. 2024}.
  \end{itemize}
  \begin{block}{Key Stellar Property}
    If a galaxy is quenched (no active star formation), it cannot host a Core-Collapse Supernova.
  \end{block}
\end{frame}

% Slide 12: Supernova Progenitors: Type Ia SNe
\begin{frame}{Supernova Progenitors: Type Ia SNe}
  \begin{itemize}
    \item \textbf{Progenitors:} Carbon-Oxygen white dwarfs in binary systems.
    \item \textbf{Explosion Trigger:} Mass transfer from a giant donor star (single-degenerate) or the merger of two white dwarfs (double-degenerate).
    \item \textbf{Delay Time Distribution (DTD):}
      Unlike massive stars, white dwarfs can take a long time to explode. The delay time $\tau$ follows a power-law distribution $\text{DTD}(\tau) \propto \tau^{-\gamma}$ (with $\gamma \approx 1.15$), spanning from $\sim 100\text{ Myr}$ to $> 10\text{ Gyr}$.
    \item \textbf{Host Environment:}
      SNe Ia occur in both active, young star-forming galaxies and old, dead elliptical galaxies. The rate depends on both the total stellar mass (older stars) and the Star Formation Rate (younger stars).
    \item \textbf{Data source:} SDSS-II Supernova Survey catalog (\textit{Lampeitl et al. 2010}).
  \end{itemize}
\end{frame}

% Slide 13: Coordinate Cross-Matching Geometry
\begin{frame}{Coordinate Cross-Matching Geometry}
  Supernovae explode within the interstellar medium of their host galaxies, often offset by several kiloparsecs from the galactic nucleus.
  \begin{itemize}
    \item On the sky, this offset corresponds to angular separations of several arcseconds.
    \item We must associate each supernova coordinate ($\text{ra}_{\text{SN}}, \text{dec}_{\text{SN}}$) with the persistent host galaxy.
    \item \textbf{Spatial Search Box:}
      We define an angular bounding box of 15 arcseconds ($0.0042^\circ$):
      $$\Delta \alpha = |\text{ra}_{\text{host}} - \text{ra}_{\text{SN}}| \le 0.0042^\circ$$
      $$\Delta \delta = |\text{dec}_{\text{host}} - \text{dec}_{\text{SN}}| \le 0.0042^\circ$$
    \item If multiple galaxies fall within this box, we calculate the exact spherical angular separation $\Delta \theta$ and select the closest:
      $$\Delta \theta = \sqrt{(\Delta \alpha \cos \text{dec}_{\text{SN}})^2 + (\Delta \delta)^2}$$
  \end{itemize}
\end{frame}

% Slide 14: SDSS Database Querying & Schema
\begin{frame}{SDSS Database Querying \& Schema}
  Once we match coordinates, we fetch the host galaxy properties by querying the SDSS DR17 spectroscopic database server.
  \begin{block}{SDSS Tables Used}
    \begin{itemize}
      \item \textbf{\texttt{PhotoObjAll} (p):} Contains photometric measurements (coordinates, magnitudes in $u,g,r,i,z$).
      \item \textbf{\texttt{SpecObj} (s):} Contains spectroscopic measurements (redshift $z$, spectral class).
      \item \textbf{\texttt{StellarMassFSPSFit} (m):} Contains stellar masses and star formation rates derived by fitting Flexible Stellar Population Synthesis (FSPS) models to the galaxy spectra.
    \end{itemize}
  \end{block}
  \begin{block}{SQL Join Condition}
    \texttt{PhotoObjAll p JOIN SpecObj s ON p.objid = s.bestObjID JOIN StellarMassFSPSFit m ON s.specObjID = m.specObjID}
  \end{block}
\end{frame}

% Slide 15: SQL WAF Bypass and Batching
\begin{frame}{SQL WAF Bypass and Batching}
  \begin{itemize}
    \item \textbf{Web Application Firewall (WAF) Block:}
      The SDSS spectroscopic server is protected by a firewall. It rejects standard multi-line SQL queries as a security precaution, returning a \texttt{403 Forbidden} error.
    \item \textbf{Bypass Implementation:}
      We must programmatically sanitize and flatten our SQL query into a single-line string, stripping out newlines (\texttt{\textbackslash n}) and excess spacing:
      \begin{itemize}
        \item \textit{Before:} \texttt{SELECT ...\textbackslash n FROM ...\textbackslash n WHERE ...}
        \item \textit{After:} \texttt{SELECT ... FROM ... WHERE ...}
      \end{itemize}
    \item \textbf{Batching:}
      To avoid \texttt{503 Gateway Timeouts} when querying many coordinates, we split our list of 50 SNe Ia coordinates into batches of 10 and query them sequentially.
    \item We implement rate-limiting (pause 1s between queries) and connection retry loops.
  \end{itemize}
\end{frame}

% Slide 16: Section 3 Divider
\begin{frame}[plain]
  \vfill
  \centering
  \begin{beamercolorbox}[sep=8pt,center,shadow=true,rounded=true]{palette primary}
    \usebeamerfont{title}\textbf{Section 3: Statistical Testing \& Results (Task 9)}
  \end{beamercolorbox}
  \vfill
\end{frame}

% Slide 17: Non-Parametric Statistics on Observational Catalogs
\begin{frame}{Non-Parametric Statistics on Observational Catalogs}
  Why can't we use standard student's t-tests or ANOVA to compare our supernova hosts with background galaxies?
  \begin{itemize}
    \item \alert{The Assumption Violation:}
      Parametric tests assume the underlying distributions are Normal (Gaussian).
    \item Galaxy stellar masses follow a Schechter function (power-law at low masses, exponential cutoff at high masses).
    \item Galaxy star-formation rates are highly bimodal (star-forming main sequence vs. quenched red sequence).
    \item \textbf{The Solution:} We must use non-parametric tests like the 1D/2D Kolmogorov-Smirnov (K-S) and Anderson-Darling (A-D) tests, which make no assumptions about the shape of the distribution.
  \end{itemize}
\end{frame}

% Slide 18: Interpreting the Mass-sSFR Plane Results
\begin{frame}{Interpreting the Mass--sSFR Plane Results}
  Executing our verification code yields the following results:
  \begin{table}
    \centering
    \begin{tabular}{lll}
      \toprule
      Test / Variable & CC-SNe Hosts vs. Background & Type Ia Hosts vs. Background \\
      \midrule
      \textbf{1D K-S Stellar Mass} & $p \approx 0.08$ (Not Significant) & $p \approx 10^{-3}$ (Highly Significant) \\
      \textbf{1D K-S sSFR}         & $p \approx 10^{-4}$ (Highly Significant) & $p \approx 0.05$ (Marginal) \\
      \textbf{2D K-S Joint Plane}  & \alert{$p \approx 10^{-5}$} (Very Significant) & \alert{$p \approx 10^{-6}$} (Extremely Significant) \\
      \bottomrule
    \end{tabular}
  \end{table}
  \begin{block}{Physical Explanation}
    \begin{itemize}
      \item \textbf{CC-SNe:} Have masses similar to average galaxies, but have much higher star-formation rates.
      \item \textbf{Type Ia:} Reside in more massive galaxies due to the rate scaling with mass, but have a mixed star-formation profile.
      \item \textbf{2D K-S Power:} By evaluating the joint Mass-sSFR plane, the 2D test captures the correlation structure (SFMS vs. Red Sequence) and reaches high statistical significance.
    \end{itemize}
  \end{block}
\end{frame}

% Slide 19: Section 4 Divider
\begin{frame}[plain]
  \vfill
  \centering
  \begin{beamercolorbox}[sep=8pt,center,shadow=true,rounded=true]{palette primary}
    \usebeamerfont{title}\textbf{Section 4: Coding Tasks \& Exercise Details}
  \end{beamercolorbox}
  \vfill
\end{frame}

% Slide 20a: Task 7 - Student Code Exercise: Mass-to-Light Ratio
\begin{frame}[fragile]{Task 7: Student Code Exercise - Mass-to-Light Ratio}
  In \texttt{stellar\_trace\_assignment2.ipynb}, you must fill in the selection effects model:
  \begin{block}{Missing Code 1: Mass-to-Light Ratio}
\begin{lstlisting}[language=Python,basicstyle=\tiny\ttfamily,keywordstyle=\color{blue}]
# Compute rest-frame g-r color with scatter
color_gr = np.clip(0.5 - 0.25 * (ssfr + 10.0), -0.1, 1.0)
color_gr += np.random.normal(0, 0.08, size=len(color_gr))

# Calculate mass-to-light ratios (Bell & de Jong 2001)
log_ml_g = 1.5 * color_gr - 0.7
log_ml_r = log_ml_g - (color_gr / 2.5)
\end{lstlisting}
  \end{block}
\end{frame}

% Slide 20b: Task 7 - Student Code Exercise: Apparent Magnitude
\begin{frame}[fragile]{Task 7: Student Code Exercise - Apparent Magnitude}
  In \texttt{stellar\_trace\_assignment2.ipynb}, compute the absolute and apparent magnitudes:
  \begin{block}{Missing Code 2: Apparent Magnitude}
\begin{lstlisting}[language=Python,basicstyle=\tiny\ttfamily,keywordstyle=\color{blue}]
# Solar absolute magnitude in r-band
m_sun_r = 4.65
absolute_mag_r = m_sun_r - 2.5 * (mass - log_ml_r)

# Apparent magnitude using distance modulus
apparent_mag_r = absolute_mag_r + 5.0 * np.log10(dl) + 25.0
\end{lstlisting}
  \end{block}
\end{frame}

% Slide 21: Task 8 - Coordinate Cross-Matching Exercise
\begin{frame}[fragile]{Task 8: Coordinate Cross-Matching Exercise}
  You must implement the spatial match function to link supernovae to their hosts:
  \begin{block}{Missing Code: Bounding Box and Distance}
\begin{lstlisting}[language=Python,basicstyle=\tiny\ttfamily,keywordstyle=\color{blue}]
# Define box limit (15 arcseconds = 0.00417 degrees)
box_limit = 15.0 / 3600.0

# Mask galaxies inside the bounding box
box_mask = (np.abs(galaxy_ra - sn_ra) <= box_limit) & \
           (np.abs(galaxy_dec - sn_dec) <= box_limit)

# Calculate spherical angular distance if matches found
if np.any(box_mask):
    matched_idx = np.where(box_mask)[0]
    ra_diff = galaxy_ra[box_mask] - sn_ra
    dec_diff = galaxy_dec[box_mask] - sn_dec
    
    # Spherical approximation
    dist = np.sqrt((ra_diff * np.cos(np.radians(sn_dec)))**2 + dec_diff**2)
    best_match = matched_idx[np.argmin(dist)]
\end{lstlisting}
  \end{block}
\end{frame}

% Slide 22: Task 8 - SQL Query Sanitizer Exercise
\begin{frame}[fragile]{Task 8: SQL Query Sanitizer Exercise}
  To bypass the SDSS Web Application Firewall (WAF), you must flatten your query:
  \begin{block}{Missing Code: SQL Sanitizer}
\begin{lstlisting}[language=Python,basicstyle=\tiny\ttfamily,keywordstyle=\color{blue}]
def sanitize_sql(query_string):
    # Strip comments starting with --
    lines = [line.split('--')[0] for line in query_string.split('\n')]
    
    # Rejoin with spaces
    flat_query = " ".join(lines)
    
    # Remove multiple consecutive spaces
    flat_query = " ".join(flat_query.split())
    return flat_query
\end{lstlisting}
  \end{block}
  \begin{exampleblock}{Mentor Tip}
    Explain to students that the WAF parses incoming strings for signature patterns. A clean, single-line string with standard syntax avoids triggers.
  \end{exampleblock}
\end{frame}

% Slide 23: Task 9 - 2D K-S Test Implementation Exercise
\begin{frame}[fragile]{Task 9: 2D K-S Test Implementation Exercise}
  Students must fill in the 2D K-S test quadrants calculation:
  \begin{block}{Missing Code: Quadrants Split}
\begin{lstlisting}[language=Python,basicstyle=\tiny\ttfamily,keywordstyle=\color{blue}]
# Evaluated at reference point (x0, y0)
q1_1 = np.sum((x1 >= x0) & (y1 >= y0)) / n1  # Top-Right
q2_1 = np.sum((x1 < x0)  & (y1 >= y0)) / n1  # Top-Left
q3_1 = np.sum((x1 < x0)  & (y1 < y0))  / n1  # Bottom-Left
q4_1 = np.sum((x1 >= x0) & (y1 < y0))  / n1  # Bottom-Right

q1_2 = np.sum((x2 >= x0) & (y2 >= y0)) / n2
q2_2 = np.sum((x2 < x0)  & (y2 >= y0)) / n2
q3_2 = np.sum((x2 < x0)  & (y2 < y0))  / n2
q4_2 = np.sum((x2 >= x0) & (y2 < y0))  / n2

# Maximum difference across the quadrants
d_local = max(abs(q1_1 - q1_2), abs(q2_1 - q2_2),
              abs(q3_1 - q3_2), abs(q4_1 - q4_2))
\end{lstlisting}
  \end{block}
\end{frame}

% Slide 24: Common Student Gotchas & Troubleshooting
\begin{frame}{Common Student Gotchas \& Troubleshooting}
  As a mentor, guide students through these common bugs:
  \begin{itemize}
    \item \textbf{Astroquery Server Down:}
      Sometimes SDSS CasJobs goes offline for maintenance. Implement local data fallbacks (`sharma\_hosts.dat`) so students can complete the lab.
    \item \textbf{Astropy Cosmology Warnings:}
      Interpolating outside the redshift grid boundaries ($z > 1.5$) throws errors. Clip redshift values to $[0.0, 1.5]$ before interpolating.
    \item \textbf{SciPy A-D Warning:}
      \texttt{scipy.stats.anderson\_ksamp} throws a warning when the p-value is outside the interpolation bounds ($p < 0.001$ or $p > 0.25$). Inform students this is normal and not a code bug.
    \item \textbf{Peacock 2D Runtime:}
      A naive Peacock 2D K-S test scales as $O(N^3)$. Warn students not to run the test on the entire background catalog ($N \sim 40,000$), but only on a sub-sampled background of $N \sim 500$!
  \end{itemize}
\end{frame}

% Slide 25: Meet 1 Summary \& Homework
\begin{frame}{Meet 1 Summary \& Homework}
  \begin{block}{What We Covered Today}
    \begin{itemize}
      \item How telescopes select light (apparent magnitude, flux limit, Malmquist Bias).
      \item Progenitor physics (Core-Collapse SNe require active star formation; Type Ia occur in both active and passive galaxies).
      \item Coordinate cross-matching using spatial search boxes.
      \item Querying SDSS and sanitizing queries to bypass firewalls.
      \item 1D and 2D non-parametric tests on host catalogs.
    \end{itemize}
  \end{block}
  \begin{description}
    \item[Homework:] Open \texttt{stellar\_trace\_assignment2.ipynb} and complete the missing codes for Tasks 7--9. Ensure your plots are saved in the \texttt{plots/} folder.
    \item[Meet 2 Preview:] Next meet, we will address class imbalances using SMOTE under strict train-test separation to prevent data leakage, and train our first Multi-Layer Perceptron classifier.
  \end{description}
\end{frame}

\end{document}
